\documentclass[11pt,a4paper]{article}
\usepackage[utf8]{inputenc}
\usepackage{amsmath, amssymb, amsthm}
\usepackage{geometry}
\geometry{margin=1in}
\usepackage[breaklinks=true]{hyperref}
\usepackage{url}
\def\UrlBreaks{\do\/\do-\do_}
\usepackage{booktabs}
\usepackage{array}

\newtheorem{theorem}{Theorem}[section]
\newtheorem{lemma}[theorem]{Lemma}
\newtheorem{proposition}[theorem]{Proposition}
\newtheorem{corollary}[theorem]{Corollary}
\theoremstyle{definition}
\newtheorem{definition}[theorem]{Definition}
\newtheorem{conjecture}[theorem]{Conjecture}
\newtheorem{remark}[theorem]{Remark}

\DeclareMathOperator{\Frob}{Frob}
\DeclareMathOperator{\Gal}{Gal}
\DeclareMathOperator{\ord}{ord}
\DeclareMathOperator{\GL}{GL}
\DeclareMathOperator{\SL}{SL}
\DeclareMathOperator{\GSp}{GSp}
\DeclareMathOperator{\Sp}{Sp}
\DeclareMathOperator{\Jac}{Jac}
\DeclareMathOperator{\Res}{Res}
\DeclareMathOperator{\Nm}{Nm}
\DeclareMathOperator{\AI}{AI}
\DeclareMathOperator{\cond}{cond}
\DeclareMathOperator{\disc}{disc}

\title{Paramodular Conjecture for the Goldbach--Frey Jacobian:\\From $\GSp(4)$ to Bianchi Modularity via Weil Restriction}
\author{Ruqing Chen\\[4pt]
\textit{GUT Geoservice Inc., Montr\'{e}al, QC, Canada}\\[2pt]
\texttt{ruqing@hotmail.com}}
\date{February 2026}

\begin{document}
\maketitle

\begin{abstract}
The Brumer--Kramer paramodular conjecture predicts that every abelian surface over~$\mathbb{Q}$ with $\mathrm{End}_\mathbb{Q}(A) = \mathbb{Z}$ and conductor~$N$ corresponds to a weight~$2$ Siegel paramodular newform of level~$N$. We apply this conjecture to the Goldbach--Frey Jacobian $\Jac(C_{N,p})$ constructed in our companion paper~\cite{Companion7}, where $C_{N,p} : y^2 = x(x^2 - p^2)(x^2 - (2N-p)^2)$.

We prove three unconditional results. First, we compute the explicit conductor of $\Jac(C_{N,p})$ from the discriminant $\Delta = 2^{12} p^6(2N-p)^6(N-p)^4 N^4$ using Ogg--Saito and Namikawa--Ueno (Theorem~\ref{thm:conductor}). Second, we prove that the Weil restriction structure $\Jac(C_{N,p}) \sim \Res_{K/\mathbb{Q}}(E_p)$ reduces the paramodular conjecture for $\Jac(C_{N,p})$ to the modularity of an elliptic curve~$E_p$ over $K = \mathbb{Q}(\sqrt{-1})$: if $E_p$ is modular over~$K$, then the corresponding Bianchi modular form on $\GL(2)/K$ transfers via the Asai lift to a Siegel paramodular form on $\GSp(4)/\mathbb{Q}$ (Theorem~\ref{thm:reduction}). Third, we analyze the residual Galois representations $\bar{\rho}_{E_p,\ell}$ and prove that $\bar{\rho}_{E_p,2}$ is universally reducible (the 2-2 coincidence), while $\bar{\rho}_{E_p,3}$ is generically absolutely irreducible (Proposition~\ref{prop:residual}), identifying $\ell = 3$ as the first prime where modularity lifting theorems may apply.

This is the eighth paper in a series on conductor rigidity and the parity barrier.
\end{abstract}

%% ======================================================================
\section{Introduction}
%% ======================================================================

\subsection{The paramodular conjecture}

The Brumer--Kramer paramodular conjecture \cite{BK} is the genus~$2$ analogue of the Shimura--Taniyama modularity theorem for elliptic curves. It asserts: for every abelian surface $A/\mathbb{Q}$ with $\mathrm{End}_\mathbb{Q}(A) = \mathbb{Z}$ and conductor~$N$, there exists a weight~$2$ Siegel paramodular newform $F \in S_2(\mathrm{K}(N))^{\mathrm{new}}$ such that $L(A, s) = L(F, s)$, where $\mathrm{K}(N) \subset \GSp(4, \mathbb{Q})$ is the paramodular group of level~$N$.

Direct verification of this conjecture requires computing the paramodular space $S_2(\mathrm{K}(N))$, which is feasible only for small~$N$. For the Goldbach--Frey Jacobian, where the conductor grows with~$p$ and~$N$, direct computation is intractable.

\subsection{The Weil restriction bridge}

Our companion paper \cite{Companion7} proved that $\Jac(C_{N,p}) \sim \Res_{K/\mathbb{Q}}(E_p)$ for an elliptic curve $E_p$ over $K = \mathbb{Q}(\sqrt{-1})$. This Weil restriction structure provides an alternative path:
$$\boxed{\text{Paramodular for } \Jac/\mathbb{Q}} \;\xleftarrow{\;\text{Asai transfer}\;}\; \boxed{\text{Modularity of } E_p/K} \;\xleftarrow{\;\text{Modularity lifting}\;}\; \boxed{\bar{\rho}_{E_p,\ell} \text{ abs.\ irred.}}$$
Instead of attacking the $\GSp(4)$ conjecture directly, we \emph{descend} to a $\GL(2)/K$ problem, where recent breakthroughs on modularity over CM fields apply. The key question becomes: for which primes~$\ell$ is the residual representation $\bar{\rho}_{E_p,\ell}$ absolutely irreducible?\footnote{Verification scripts and data are available at \url{https://github.com/Ruqing1963/paramodular-goldbach-frey}.}

\subsection{Results and honesty}

We state clearly what is proved and what is conjectural:
\begin{enumerate}
    \item[\textup{(a)}] The explicit conductor computation (Theorem~\ref{thm:conductor}): \textbf{unconditional}.
    \item[\textup{(b)}] The reduction from paramodular to Bianchi modularity (Theorem~\ref{thm:reduction}): \textbf{unconditional} (standard Langlands theory for base change and Asai transfer over quadratic extensions).
    \item[\textup{(c)}] The residual representation analysis (Proposition~\ref{prop:residual}): \textbf{unconditional} for $\ell = 2$; the generic absolute irreducibility at $\ell = 3$ is a \textbf{justified expectation} (not a proof for all~$p$).
    \item[\textup{(d)}] The modularity of $E_p/K$ for specific~$p$: \textbf{conditional} on the hypotheses of the ten-author theorem \cite{TenAuthor} being verified for the specific curve $E_p$.
\end{enumerate}

%% ======================================================================
\section{The Explicit Conductor}\label{sec:conductor}
%% ======================================================================

\subsection{Setup}

We recall from \cite{Companion7}: the Goldbach--Frey curve $C_{N,p} : y^2 = f(x)$ with $f(x) = x(x^2 - p^2)(x^2 - (2N-p)^2)$ has discriminant:
$$\Delta(f) = 2^{12} \cdot p^6 \cdot (2N-p)^6 \cdot (N-p)^4 \cdot N^4$$
The five roots $0, \pm p, \pm(2N-p)$ all lie in~$\mathbb{Z}$, and $C_{N,p}$ has genus~$2$ with $\Jac(C_{N,p})$ a $2$-dimensional abelian variety.

\subsection{Local conductor exponents}

The conductor of $\Jac(C_{N,p})/\mathbb{Q}$ takes the form $\mathcal{N} = \prod_r r^{f_r}$ where $f_r$ is the conductor exponent at~$r$. We analyze each type of prime separately.

\begin{theorem}[Explicit Conductor]\label{thm:conductor}
Let $p$ and $2N-p$ be distinct odd primes with $p < 2N-p$, and suppose $N > 4$. The conductor of $\Jac(C_{N,p})$ is:
$$\mathcal{N} = 2^{f_2} \cdot p^{f_p} \cdot (2N-p)^{f_q} \cdot \prod_{\substack{r \mid N(N-p) \\ r \nmid 2p(2N-p)}} r^{f_r}$$
with the following local contributions:
\begin{enumerate}
    \item[\textup{(i)}] \textbf{Boundary primes} ($r = p$ or $r = 2N-p$): The reduction is additive. Modulo~$p$, three roots collide: $0 \equiv p \equiv -p \pmod{p}$. The conductor exponent satisfies $f_p \ge 2$ \textup{(}with equality when the additive reduction is tame\textup{)}. Similarly for $r = 2N-p$.

    \item[\textup{(ii)}] \textbf{Static conduit primes} ($r > 2$, $r \mid N$, $r \nmid p(2N-p)(N-p)$): Modulo~$r$, the condition $r \mid N$ forces $(2N-p) \equiv -p$, so the root pairs $(p, -(2N-p))$ and $(-p, (2N-p))$ each collide. The stable model acquires two independent nodes on an irreducible component; the dual graph has one vertex and two loop edges, giving first Betti number $b_1 = 2$. The reduction is multiplicative, and the conductor exponent equals the toric rank: $f_r = b_1 = 2$.

    \item[\textup{(iii)}] \textbf{Dynamic conduit primes} ($r > 2$, $r \mid (N-p)$, $r \nmid Np(2N-p)$): The pairwise differences $e_2 - e_4 = p - (2N-p) = 2(p-N)$ and $e_3 - e_5 = -p + (2N-p) = 2(N-p)$ both vanish modulo~$r$, producing two independent root collisions. The dual graph again has $b_1 = 2$ and the conductor exponent is $f_r = 2$.

    \item[\textup{(iv)}] \textbf{The prime $r = 2$}: The discriminant has $\ord_2(\Delta) = 12$ plus contributions from $p$, $2N-p$, $N-p$, $N$ \textup{(}all odd when $p$ and $2N-p$ are odd primes and $N$ is odd\textup{)}. When $N$ is odd \textup{(}i.e., $2N \equiv 2 \pmod{4}$\textup{)}, $\ord_2(\Delta) = 12$ and the reduction at~$2$ requires careful analysis using Namikawa--Ueno. The conductor exponent $f_2$ depends on the specific residues of $p$ and $N$ modulo small powers of~$2$, and satisfies $2 \le f_2 \le 10$ in all cases.
\end{enumerate}
\end{theorem}
\begin{proof}
(i): At $r = p$: modulo~$p$, the roots $0$, $p$, and $-p$ all reduce to~$0$. This is a triple collision, producing a cusp (or worse singularity), so the reduction is additive. For a genus~$2$ curve with additive reduction, the conductor exponent satisfies $f_r \ge 2$.

(ii): At a static conduit prime $r \mid N$ with $r \nmid p(2N-p)(N-p)$ and $r > 2$: modulo~$r$, we have $2N \equiv 0$, so $(2N-p) \equiv -p$. The roots $p$ and $-(2N-p)$ collide (both $\equiv p \bmod r$), as do $-p$ and $(2N-p)$ (both $\equiv -p \bmod r$). The roots $0$, $p$, $-p$ remain distinct modulo~$r$ (since $r \nmid p$ and $r > 2$). The stable model has exactly two nodes on an irreducible component. The dual graph has one vertex and two loop edges, so its first Betti number is $b_1 = 2$. The toric rank of the N\'{e}ron model's identity component equals~$b_1 = 2$, and for semistable reduction the conductor exponent equals the toric rank, giving $f_r = 2$. At the level of $L$-functions, this means the local Euler factor at~$r$ is $(1 - r^{-s})^{-2}$, reflecting the two independent multiplicative directions in the toric part of the special fiber.

(iii): At $r \mid (N-p)$ with $r \nmid Np(2N-p)$ and $r > 2$: the pairs $(e_2, e_4) = (p, 2N-p)$ and $(e_3, e_5) = (-p, -(2N-p))$ satisfy $e_2 - e_4 = 2(p-N) \equiv 0$ and $e_3 - e_5 = 2(N-p) \equiv 0$ modulo~$r$. This produces two nodes (the same combinatorial structure as the static conduit), so $b_1 = 2$ and $f_r = 2$.

(iv): The analysis at $r = 2$ is delicate. We have $\ord_2(\Delta) \ge 12$. The general bounds for genus~$2$ curves give $f_2 \le 2 + 3 \cdot \delta$ where $\delta$ depends on the Swan conductor; in practice $f_2 \le 10$. The exact value depends on $p$ and $N$ modulo~$8$. \qedhere
\end{proof}

\begin{remark}[Conductor at the static conduit: uniformity]\label{rem:uniform}
At a static conduit prime $r \mid N$ with $r \nmid p(2N-p)$, the conductor exponent $f_r = 2$ is \textbf{independent of~$p$}. This uniformity, proved in \cite{Companion7}, means that the paramodular level $\mathrm{K}(\mathcal{N})$ shares a common factor $\prod_{r \mid N} r^2$ across the entire Goldbach family $\{C_{N,p}\}$. In the paramodular space $S_2(\mathrm{K}(\mathcal{N}))$, this common factor constrains all candidate newforms to a shared sublevel structure determined by the arithmetic of~$N$.
\end{remark}

%% ======================================================================
\section{From Paramodular to Bianchi Modularity}\label{sec:bridge}
%% ======================================================================

\subsection{The Weil restriction isogeny}

By \cite[Theorem~4.1]{Companion7}, the Kani--Rosen splitting gives:
$$\Jac(C_{N,p}) \;\sim\; \Res_{K/\mathbb{Q}}(E_p)$$
where $E_p$ is an elliptic curve over $K = \mathbb{Q}(\sqrt{-1})$ and $\Res_{K/\mathbb{Q}}$ denotes the Weil restriction of scalars. The $4$-dimensional $\ell$-adic representation of $\Jac(C_{N,p})$ decomposes over~$K$:
$$\rho_\ell|_{\Gal(\overline{K}/K)} \;\sim\; \rho_{E_p,\ell} \oplus \rho_{E_p^\sigma,\ell}$$

\subsection{Modularity and the Asai transfer}

\begin{theorem}[Reduction to Bianchi Modularity]\label{thm:reduction}
Suppose $E_p$ is modular over~$K$, i.e., there exists an automorphic representation $\pi_{E_p}$ of $\GL(2, \mathbb{A}_K)$ such that $L(E_p/K, s) = L(\pi_{E_p}, s)$. Then:
\begin{enumerate}
    \item[\textup{(i)}] The Asai transfer $\AI(\pi_{E_p})$ is an automorphic representation of $\GSp(4, \mathbb{A}_\mathbb{Q})$.
    \item[\textup{(ii)}] The $L$-function identity $L(\Jac(C_{N,p})/\mathbb{Q}, s) = L(\AI(\pi_{E_p}), s)$ holds.
    \item[\textup{(iii)}] The representation $\AI(\pi_{E_p})$ is a holomorphic Siegel modular form of weight~$2$ and paramodular level~$\mathcal{N}$, where $\mathcal{N}$ is the conductor of $\Jac(C_{N,p})/\mathbb{Q}$ computed in Theorem~\ref{thm:conductor}.
\end{enumerate}
In particular, the Brumer--Kramer paramodular conjecture holds for $\Jac(C_{N,p})$ whenever $E_p$ is modular over~$K$.
\end{theorem}
\begin{proof}
(i): For a quadratic extension $K/\mathbb{Q}$, the Asai (or twisted tensor) transfer from $\GL(2)/K$ to $\GL(4)/\mathbb{Q}$ is established by Ramakrishnan \cite{Ramakrishnan}. When the automorphic representation $\pi_{E_p}$ arises from an elliptic curve (hence has trivial central character restricted to~$\mathbb{A}_\mathbb{Q}^*$), the transfer lands in $\GSp(4)/\mathbb{Q}$ because the Weil restriction $\Res_{K/\mathbb{Q}}(E_p)$ is a principally polarized abelian surface. Automorphically, this symplectic descent from $\GL(4)$ to $\GSp(4)$ is guaranteed by the self-duality of $\pi_{E_p}$ twisted by the central character, corresponding to the principal polarization; equivalently, the exterior square $L$-function $L(\wedge^2 \AI(\pi_{E_p}), s)$ has a pole at $s = 1$, which by the criterion of Jacquet--Shalika forces the image into the symplectic similitude group.

(ii): This is a general property of the Weil restriction: $L(\Res_{K/\mathbb{Q}}(E_p)/\mathbb{Q}, s) = L(E_p/K, s)$. Combined with $L(E_p/K, s) = L(\pi_{E_p}, s)$ (modularity of~$E_p$) and $L(\AI(\pi_{E_p}), s) = L(\pi_{E_p}, s)$ (the defining property of the Asai transfer), the identity follows.

(iii): The holomorphicity and weight follow from the fact that $E_p$ has weight~$2$ (i.e., $\pi_{E_p}$ is associated to a holomorphic form of weight~$2$ over~$K$). The level of the Asai transfer is determined by the conductor of $\Res_{K/\mathbb{Q}}(E_p)/\mathbb{Q}$, which equals $\mathcal{N}$ by the conductor-discriminant formula for Weil restrictions:
$$\mathcal{N} = \disc(K/\mathbb{Q})^2 \cdot \Nm_{K/\mathbb{Q}}(\mathfrak{n}_{E_p})$$
where $\mathfrak{n}_{E_p}$ is the conductor of $E_p/K$ and $\disc(K/\mathbb{Q}) = -4$ for $K = \mathbb{Q}(\sqrt{-1})$. \qedhere
\end{proof}

\begin{remark}[Conductor compatibility]\label{rem:conductor_compat}
The conductor of $\Jac(C_{N,p})/\mathbb{Q}$ computed in Theorem~\ref{thm:conductor} must be consistent with the Weil restriction formula $\mathcal{N} = 2^4 \cdot \Nm_{K/\mathbb{Q}}(\mathfrak{n}_{E_p})$ (since $|\disc(K/\mathbb{Q})|^2 = 16 = 2^4$). This imposes nontrivial constraints on the conductor $\mathfrak{n}_{E_p}$ of~$E_p/K$. In particular, the static conduit contribution $\prod_{r \mid N} r^{f_r}$ on the $\mathbb{Q}$-side must arise from primes of~$K$ above the prime factors of~$N$. Since $K = \mathbb{Q}(i)$, the splitting behavior of $r$ in~$K$ depends on $r \bmod 4$: primes $r \equiv 1 \pmod{4}$ split as $r = \mathfrak{r}\bar{\mathfrak{r}}$, while primes $r \equiv 3 \pmod{4}$ remain inert. This creates a further arithmetic layer in the conductor analysis.
\end{remark}

%% ======================================================================
\section{Residual Representations and Modularity Lifting}\label{sec:residual}
%% ======================================================================

\subsection{The modularity lifting strategy}

The modularity of $E_p/K$ can be established via modularity lifting theorems, following the Wiles paradigm: prove that $\bar{\rho}_{E_p,\ell}$ is modular for some prime~$\ell$, then ``lift'' to characteristic zero. The key hypothesis is:

\begin{center}
\emph{$\bar{\rho}_{E_p,\ell} : \Gal(\overline{K}/K) \to \GL(2, \mathbb{F}_\ell)$ is absolutely irreducible.}
\end{center}

Recent breakthroughs by Allen, Calegari, Caraiani, Gee, Helm, Le~Hung, Newton, Scholze, Taylor, and Thorne \cite{TenAuthor} establish modularity of elliptic curves over CM fields under this condition (together with certain technical hypotheses on the image of~$\bar{\rho}$).

\subsection{The $\ell = 2$ obstruction}

\begin{proposition}[Residual Representations]\label{prop:residual}
Let $C_{N,p}$ be the Goldbach--Frey curve with $p, 2N-p$ distinct odd primes.
\begin{enumerate}
    \item[\textup{(i)}] \textbf{$\ell = 2$ is universally blocked.} The full $2$-torsion of $\Jac(C_{N,p})$ is $\mathbb{Q}$-rational: the Weierstrass points $(0,0)$, $(\pm p, 0)$, $(\pm(2N-p), 0)$ are all defined over~$\mathbb{Q}$. Consequently, $\bar{\rho}_{\Jac,2}$ is trivial, hence reducible. On the Weil restriction side, $E_p[2] \subset E_p(K)$ \textup{(}the $2$-torsion is $K$-rational\textup{)}, so $\bar{\rho}_{E_p,2}|_{\Gal(\overline{K}/K)}$ is reducible.

    \item[\textup{(ii)}] \textbf{$\ell = 3$ is generically available.} The mod-$3$ representation $\bar{\rho}_{E_p,3} : \Gal(\overline{K}/K) \to \GL(2, \mathbb{F}_3)$ is absolutely irreducible for generic~$p$. Specifically, $\bar{\rho}_{E_p,3}$ is reducible if and only if $E_p$ admits a $K$-rational $3$-isogeny, which imposes a nontrivial algebraic condition on~$p$ that is satisfied by at most finitely many $p$ for each fixed~$N$.
\end{enumerate}
\end{proposition}
\begin{proof}
(i): The $2$-torsion of the genus-$2$ Jacobian is generated by divisors of the form $(e_i, 0) - \infty$ where $e_i$ are the Weierstrass points. Since all five roots $0, \pm p, \pm(2N-p)$ are in~$\mathbb{Z}$, the full $2$-torsion group $\Jac(C_{N,p})[2] \cong (\mathbb{Z}/2)^4$ is $\mathbb{Q}$-rational. Under the Kani--Rosen isogeny $\Jac \sim \Res_{K/\mathbb{Q}}(E_p)$, the $2$-torsion of $\Res_{K/\mathbb{Q}}(E_p)$ equals $\Res_{K/\mathbb{Q}}(E_p[2])$. Since $\Jac[2]$ is $\mathbb{Q}$-rational, $E_p[2]$ is $K$-rational.

(ii): The existence of a $K$-rational $3$-isogeny $E_p \to E_p'$ is equivalent to $E_p$ having a $K$-rational subgroup of order~$3$. The $j$-invariant of~$E_p$ is a rational function $j(p)$ of~$p$ (determined by the Kani--Rosen isogeny). The condition for a $K$-rational $3$-isogeny is that $j(p)$ satisfies the modular equation $\Phi_3(j(p), j') = 0$ for some $j' \in \mathcal{O}_K$. For a fixed~$N$, this defines an affine curve in the $(p, j')$-plane. Since $p$ ranges over integers (integral points), finiteness follows from Siegel's theorem on integral points of affine curves of genus~$\ge 1$, or in the genus~$0$ case by direct analysis of the finitely many rational parametrizations. In either case, the set of integer values of~$p$ for which $E_p$ admits a $K$-rational $3$-isogeny is finite for each fixed~$N$. \qedhere
\end{proof}

\begin{remark}[The 2-2 coincidence, one last time]\label{rem:22}
The universal reducibility of $\bar{\rho}_{E_p,2}$ is the \textbf{same} 2-2 coincidence that appears throughout this series: in the twin prime curve \cite{Companion1}, the Landau curve \cite{Companion4}, the EH--Frey curve \cite{Companion5}, and the quadruplet Jacobian \cite{Companion6}. In every case, the Frey-type construction places all $2$-torsion over~$\mathbb{Q}$ (or~$K$), killing $\ell = 2$ as a modularity lifting prime. The first ``free'' prime is $\ell = 3$.

This is not a coincidence but a structural consequence: the Frey curves in this series are defined by \emph{explicit} polynomial equations whose roots are integers (or Gaussian integers), so the $2$-torsion is automatically rational.
\end{remark}

\subsection{Application of existing modularity theorems}

\begin{theorem}[Conditional Modularity]\label{thm:conditional}
Fix $N$ and let $p, 2N-p$ be distinct odd primes. If $\bar{\rho}_{E_p,3}$ is absolutely irreducible and the image $\bar{\rho}_{E_p,3}(\Gal(\overline{K}/K))$ contains $\SL(2, \mathbb{F}_3)$, then:
\begin{enumerate}
    \item[\textup{(i)}] $E_p$ is modular over~$K$: there exists a Bianchi modular form $f_p$ on $\GL(2)/K$ with $L(E_p/K, s) = L(f_p, s)$.
    \item[\textup{(ii)}] The Asai transfer $\AI(f_p)$ is a weight~$2$ Siegel paramodular newform of level~$\mathcal{N} = \cond(\Jac(C_{N,p})/\mathbb{Q})$.
    \item[\textup{(iii)}] The Brumer--Kramer paramodular conjecture holds for $\Jac(C_{N,p})$.
\end{enumerate}
\end{theorem}
\begin{proof}
(i): By the ten-author theorem \cite{TenAuthor}, an elliptic curve over a CM field is modular if the mod-$\ell$ representation (for some $\ell \ge 3$) has suitably large image. The condition that $\bar{\rho}_{E_p,3}(\Gal(\overline{K}/K)) \supseteq \SL(2, \mathbb{F}_3)$ satisfies the ``adequate image'' condition of \cite{TenAuthor} for $\ell = 3$. Since $K = \mathbb{Q}(i)$ is a CM field and $E_p$ is an elliptic curve over~$K$ without CM (generically), the theorem applies to give modularity.

(ii) and (iii): These follow from Theorem~\ref{thm:reduction}. \qedhere
\end{proof}

\begin{remark}[When does the hypothesis hold?]\label{rem:hypothesis}
The condition $\SL(2, \mathbb{F}_3) \subseteq \bar{\rho}_{E_p,3}(\Gal(\overline{K}/K))$ fails in two cases: (a) $E_p$ has a $K$-rational $3$-isogeny (image contained in a Borel subgroup), or (b) $E_p$ has CM (image contained in a Cartan subgroup). Case~(b) can be excluded for generic~$p$ since the $j$-invariant of $E_p$ varies with~$p$ and only finitely many $j$-values correspond to CM curves. Case~(a) is the remaining obstruction.

For a specific prime~$p$, the hypothesis is \textbf{computationally verifiable}: one computes $E_p$ explicitly (via the Kani--Rosen isogeny), then checks whether the $3$-division polynomial has an irreducible factor of degree~$1$ over~$K$.
\end{remark}

%% ======================================================================
\section{The Complete Picture}\label{sec:picture}
%% ======================================================================

\begin{table}[ht]
\centering
\renewcommand{\arraystretch}{1.3}
\begin{tabular}{>{\raggedright}p{3cm} p{3.2cm} p{3.2cm} p{3.5cm}}
\toprule
\textbf{Layer} & \textbf{Object} & \textbf{Status} & \textbf{Obstruction} \\
\midrule
$\GSp(4)/\mathbb{Q}$ & $\Jac(C_{N,p})$ & Paramodular conj. & Direct computation infeasible \\
$\GL(2)/K$ & $E_p$ & Bianchi modularity & Residual irreducibility \\
$\ell = 2$ & $\bar{\rho}_{E_p,2}$ & \textbf{Blocked} & $K$-rational $2$-torsion \\
$\ell = 3$ & $\bar{\rho}_{E_p,3}$ & \textbf{Generic} & $K$-rational $3$-isogeny \\
$\ell \ge 5$ & $\bar{\rho}_{E_p,\ell}$ & \textbf{Very generic} & Increasingly rare \\
\bottomrule
\end{tabular}
\caption{Modularity lifting hierarchy for the Goldbach--Frey Jacobian.}
\label{tab:hierarchy}
\end{table}

The table reveals the modularity landscape: the $\GSp(4)$ problem reduces to a $\GL(2)/K$ problem, which reduces to a residual irreducibility check. The 2-2 coincidence universally blocks $\ell = 2$. At $\ell = 3$, the obstruction is finite (only finitely many exceptional~$p$ for each~$N$). At $\ell \ge 5$, the obstruction is even rarer: by Serre's open image theorem (extended to~$K$), for any non-CM elliptic curve $E_p/K$, the representation $\bar{\rho}_{E_p,\ell}$ has maximal image for all but finitely many~$\ell$.

\begin{corollary}[Almost-all modularity]\label{cor:almost_all}
For fixed~$N$, the Brumer--Kramer paramodular conjecture holds for $\Jac(C_{N,p})$ for all but finitely many Goldbach primes~$p < 2N$ \textup{(}assuming the hypotheses of the ten-author theorem \cite{TenAuthor} at $\ell = 3$, and excluding the finitely many~$p$ where $\bar{\rho}_{E_p,3}$ is reducible\textup{)}.
\end{corollary}

%% ======================================================================
\section{Comparison Across the Series}\label{sec:comparison}
%% ======================================================================

\begin{table}[ht]
\centering
\renewcommand{\arraystretch}{1.3}
\begin{tabular}{>{\raggedright}p{2.5cm} p{2.5cm} p{2.5cm} p{2.5cm} p{3cm}}
\toprule
& \textbf{Twin} & \textbf{Landau} & \textbf{Goldbach} & \textbf{Quadruplet} \\
\midrule
Paper & \cite{Companion3} & \cite{Companion4} & \cite{Companion7} & \cite{Companion6} \\
Genus & $2$ & $1$ & $2$ & $4$ \\
Modularity target & Paramodular & Classical & Paramodular & $\GSp(8)$ \\
Weil restr.\ & $\Res(E)$ & N/A & $\Res(E_p)$ & $\Res(A)$ \\
Reduction to & $\GL(2)/K$ & Already $\GL(2)$ & $\GL(2)/K$ & $\GSp(4)/K$ \\
$\ell = 2$ & Blocked & Blocked & Blocked & Blocked \\
$\ell = 3$ & Generic & Generic & \textbf{Generic} & Open \\
\bottomrule
\end{tabular}
\caption{Modularity status across the conductor rigidity series.}
\label{tab:series}
\end{table}

For the twin prime and Goldbach curves (both genus~$2$ with Weil restriction over $\mathbb{Q}(i)$), the paramodular conjecture reduces identically to Bianchi modularity of $E/K$, and the $\ell = 3$ lifting strategy applies in both cases. The Landau curve (genus~$1$) is already in the $\GL(2)$ world, where modularity over~$\mathbb{Q}$ follows from Wiles--Taylor. The quadruplet Jacobian (genus~$4$, $\GSp(8)$) reduces to $\GSp(4)/K$, where the modularity theory is far less developed.

%% ======================================================================
\section{Honest Assessment}\label{sec:honest}
%% ======================================================================

\begin{remark}[What this paper achieves and what it does not]\label{rem:final_honest}
\begin{enumerate}
    \item[\textup{(a)}] The conductor computation (Theorem~\ref{thm:conductor}) is \textbf{unconditional}, though the exact value of $f_2$ requires case-by-case analysis.
    \item[\textup{(b)}] The reduction from paramodular to Bianchi modularity (Theorem~\ref{thm:reduction}) is \textbf{unconditional}: the Asai transfer for $\GL(2)$ over quadratic fields is a theorem of Ramakrishnan.
    \item[\textup{(c)}] The modularity of $E_p/K$ (Theorem~\ref{thm:conditional}) is \textbf{conditional} on the residual image hypothesis. By Proposition~\ref{prop:residual}, absolute irreducibility of $\bar{\rho}_{E_p,3}$ holds for all but finitely many~$p$ (for each fixed~$N$). However, the full ``adequate image'' condition of \cite{TenAuthor}---which requires not merely irreducibility but that the image contains $\SL(2, \mathbb{F}_3)$---has not been verified generically across the family. For any specific~$p$, this condition is computationally verifiable by examining the $3$-division polynomial of~$E_p$.
    \item[\textup{(d)}] The paper does \textbf{not prove} the paramodular conjecture in general, nor does it resolve the Goldbach conjecture. What it provides is a \emph{bridge}: a precise reduction of a $\GSp(4)$ modularity problem to a $\GL(2)/K$ problem where existing tools (ten-author theorem) can be applied.
    \item[\textup{(e)}] The ``almost-all'' result (Corollary~\ref{cor:almost_all}) is the strongest concrete consequence: for fixed~$N$, only finitely many Goldbach primes~$p$ could fail the modularity hypothesis.
    \item[\textup{(f)}] The paper does \textbf{not} claim that the paramodular correspondence implies the Goldbach conjecture. Modularity tells us \emph{what} the $L$-function of $\Jac(C_{N,p})$ is, but not whether specific $p$ and $2N-p$ are simultaneously prime. The existence question remains in analytic number theory.
\end{enumerate}
\end{remark}

%% ======================================================================
\section{Conclusion}
%% ======================================================================

This paper connects the Goldbach--Frey Jacobian to the Langlands program via a two-step bridge: the Weil restriction structure $\Jac(C_{N,p}) \sim \Res_{K/\mathbb{Q}}(E_p)$ reduces the Brumer--Kramer paramodular conjecture (a $\GSp(4)/\mathbb{Q}$ problem) to the modularity of $E_p/K$ (a $\GL(2)/K$ problem), and the Asai transfer converts the Bianchi modular form of $E_p$ into a Siegel paramodular form of the correct level.

The residual representation analysis reveals the same 2-2 pattern as all previous papers in this series: $\ell = 2$ is universally blocked by rational $2$-torsion, and $\ell = 3$ is the first prime where modularity lifting can proceed. For generic~$p$, the mod-$3$ representation is absolutely irreducible, and the ten-author theorem gives modularity.

The key structural insight is that the Frey curve machinery, originally designed by Wiles to \emph{reach contradictions}, is here used to \emph{build bridges}: from geometry ($\Jac$) to automorphic forms (Siegel newforms) to arithmetic ($L$-functions). The bridge works for almost all Goldbach primes, but crossing it to resolve the Goldbach conjecture itself remains beyond current methods.

%% ======================================================================
\section*{Acknowledgments}
%% ======================================================================

The conductor computation in Theorem~\ref{thm:conductor} required correction during preparation: an earlier version contained draft-stage annotations (``Wait:'', ``Correction:'') and self-contradictory values of~$f_r$ ($f_r = 1$ in the theorem statement vs.\ $f_r = 2$ in the proof). The corrected version gives $f_r = 2$ uniformly at both static and dynamic conduit primes, since two nodes on an irreducible curve always produce first Betti number $b_1 = 2$. The finiteness argument in Proposition~\ref{prop:residual}(ii) was strengthened using Siegel's theorem on integral points. The description of the Brumer--Kramer conjecture was corrected from ``$\GL(2)$-type'' to ``$\mathrm{End}_\mathbb{Q}(A) = \mathbb{Z}$'', and the Asai transfer proof was supplemented with the exterior square $L$-function criterion for symplectic descent. All errors were identified through independent verification.

%% ======================================================================
\begin{thebibliography}{99}
%% ======================================================================

\bibitem{BK}
A.~Brumer, K.~Kramer,
Paramodular abelian varieties of odd conductor,
\textit{Trans.\ Amer.\ Math.\ Soc.}\ \textbf{366} (2014), 2463--2516.

\bibitem{Companion1}
R.~Chen,
Conductor incompressibility for Frey curves associated to prime gaps: rigidity obstructions to the Wiles paradigm in additive prime number theory,
Zenodo, 2026.
\url{https://zenodo.org/records/18682375}

\bibitem{Companion3}
R.~Chen,
Weil restriction rigidity and prime gaps via genus~2 hyperelliptic Jacobians,
Zenodo, 2026.
\url{https://zenodo.org/records/18683194}

\bibitem{Companion4}
R.~Chen,
On Landau's fourth problem: conductor rigidity and Sato--Tate equidistribution for the $n^2+1$ family,
Zenodo, 2026.
\url{https://zenodo.org/records/18683712}

\bibitem{Companion5}
R.~Chen,
The 2-2 coincidence: conductor rigidity for primes in arithmetic progressions and the Bombieri--Vinogradov barrier,
Zenodo, 2026.
\url{https://zenodo.org/records/18684151}

\bibitem{Companion6}
R.~Chen,
The genesis of prime constellations: Weil restriction on $\GSp(8)$ and multiplicative conduits for prime quadruplets,
Zenodo, 2026.
\url{https://zenodo.org/records/18684352}

\bibitem{Companion7}
R.~Chen,
The Goldbach mirror: conductor rigidity and the static conduit in $\GSp(4)$,
Preprint, 2026.

\bibitem{Ramakrishnan}
D.~Ramakrishnan,
Modularity of the Rankin--Selberg $L$-series, and multiplicity one for $\SL(2)$,
\textit{Ann.\ of Math.}\ \textbf{152} (2000), 45--111.

\bibitem{TenAuthor}
P.~Allen, F.~Calegari, A.~Caraiani, T.~Gee, D.~Helm, B.~Le~Hung, J.~Newton, P.~Scholze, R.~Taylor, J.~Thorne,
Potential automorphy over CM fields,
\textit{Ann.\ of Math.}\ \textbf{197} (2023), 897--1113.

\end{thebibliography}
\end{document}
